\SourceLeaf{278}
\par\indent Die Zahlenwerthe in unserm Beispiele (für den briggischen Logarithmen, und \(q\) in Graden ausgedrückt) sind
\[
\begin{aligned}
\log m={}&-0{,}0049796163\ 94(\frac{q}{100})^3\\
&-0{,}0016150307\ 6(\frac{q}{100})^4\\
&-0{,}0023973954(\frac{q}{100})^5\\
&-0{,}0125671(\frac{q}{100})^6
\end{aligned}
\]
\Article{10}
\indent Bei einer weitumfassenden Vermessung, wo die Übertragung vom Sphäroid auf die Kugel oder umgekehrt für sehr viele Punkte vorkommt, wird man, anstatt jedesmal auf die Formeln zurückzukommen, lieber ein für allemal eine ausgedehnte Tafel berechnen. Der Gebrauch einer solchen Tafel wird aber bequemer sein, wenn man ihr die Breite auf der Kugel \(Q+q\) zum Argument gibt, als wenn man die Breite auf dem Sphäroid dazu wählen wollte, indem der Übergang von ersterer auf die andere viel häufiger erfordert wird, als der umgekehrte. Für jeden Rechnungserfahrnen wird übrigens die Bemerkung überflüssig sein, dass man behuf Construction einer solchen Tafel nur eine mässige Anzahl von Gliedern direct berechnet, aus denen die übrigen mit eben so grosser Schärfe und sehr geringer Mühe durch ein angemessenes Interpolationsverfahren bestimmt werden. Es werden also dafür die im 8. und 9. Artikel mitgetheilten Reihen zur Anwendung kommen, und gerade deswegen ist es vortheilhaft, dass nicht \(P\), sondern \(Q\) eine runde Zahl sei.

Ich füge am Schluss dieser Abhandlung eine solche Tafel bei, welcher der Normalwerth \(Q=52^0 40'\) (wie dem bisher betrachteten Beispiele) zum Grunde liegt, und die durch zwölf Grade, von \(46^0 40'\) bis \(58^0 40'\), für alle Werthe des Arguments \(Q+q\) von Minute zu Minute fortschreitet. Sie gibt den zugehörigen Werth von \(P+p\) auf fünf Decimalen der Secunde genau; ferner den briggischen Logarithmen von \(m\) auf zehn Stellen, nemlich in Einheiten der zehnten Decimale; endlich auch noch, in Secunden ausgedrückt, den Werth von \(-\tfrac{\mathrm{d}m}{2m\mathrm{d}q}\); der Gebrauch dieser letzten Columne wird weiter unten erklärt werden. Ich habe die Tafel deshalb mit so vielen Decimalen gegeben, damit sie auch für die allerschärfste Berechnung einer trigonometrischen Vermessung, nemlich für eine
